\chapter{Design and Implementation}

This chapter details the system architecture, design specifications, and implementation details of our platform. It describes the data preprocessing pipeline, feature engineering process, hybrid model components, and the integration of the REST \gls{api} and interactive dashboard.

\section{System Overview}

FraudNet is structured as a modular pipeline that automates the transition from raw transactional records to real-time model scoring and visual explanations. The system comprises three main blocks: the data preparation block (preprocessing, feature engineering, and sequence generation), the machine learning block (model training, hyperparameter optimization, and hybrid ensemble blending), and the application block (FastAPI endpoint and Streamlit dashboard).

\section{System Requirement Specifications}

\subsection{Software Requirements}
The software requirements for developing and running the FraudNet platform are:
\begin{itemize}
    \item \textbf{Programming Language}: Python $\ge$ 3.10
    \item \textbf{Data Processing Libraries}: Pandas $\ge$ 2.0.0, NumPy $\ge$ 1.24.0, PyArrow $\ge$ 14.0.0
    \item \textbf{Machine Learning Libraries}: Scikit-learn $\ge$ 1.3.0, XGBoost $\ge$ 2.0.0, LightGBM $\ge$ 4.0.0, Imbalanced-learn $\ge$ 0.12.0
    \item \textbf{Explainability Tool}: SHAP $\ge$ 0.45.0
    \item \textbf{API \& UI Frameworks}: FastAPI $\ge$ 0.115.0, Uvicorn $\ge$ 0.30.0, Streamlit $\ge$ 1.36.0
    \item \textbf{Testing Framework}: Pytest $\ge$ 7.4.0
\end{itemize}

\subsection{Functional Requirements}
\begin{itemize}
    \item \textbf{Data Preprocessing}: Clean, deduplicate, scale, and save transactional data automatically.
    \item \textbf{Model Training \& Tuning}: Train seven distinct ML models and perform hyperparameter searches.
    \item \textbf{Real-Time Scoring}: Serve transaction scores via an HTTP REST endpoint.
    \item \textbf{Interpretability}: Generate SHAP force plots explaining features dynamically.
\end{itemize}

\subsection{Non-Functional Requirements}
\begin{itemize}
    \item \textbf{Low Latency}: Score transactions in under 50 milliseconds via the API.
    \item \textbf{Leakage Protection}: Prevent target leakage by performing data oversampling strictly inside cross-validation splits.
    \item \textbf{Scalability}: Support batch dataset uploads and scoring in the dashboard.
\end{itemize}

\section{System Design}

\subsection{Architectural Design}
The system architecture of FraudNet is separated into training and inference phases.
The training pipeline ingests raw Parquet datasets, splits them using group-aware strategies, constructs sequential inputs, trains components, and serializes trained estimators to disk. 
The inference pipeline loads these serialized estimators to instantiate the `FraudNetScoringService` which routes predictions using a 5-agent conceptual trace:

\begin{enumerate}
    \item \textbf{Packet Agent}: Receives and validates raw inputs (either 29 raw features or 36 pre-engineered features).
    \item \textbf{Flow Agent}: Applies scaling, logarithmic, and frequency feature engineering to raw inputs.
    \item \textbf{Behavior Agent}: Calls the \gls{rf} component and sequence-based \gls{rnn} component to retrieve component scores.
    \item \textbf{Correlation Agent}: Blends component scores using optimized weights.
    \item \textbf{Response Agent}: Compiles the final classification decision and \gls{shap} interpretability outputs.
\end{enumerate}

\subsection{Module or Component Design}
The software modules are organized as follows:
\begin{itemize}
    \item \textbf{src/preprocessing.py}: Cleans raw transaction files, checks for missing data, and performs deduplication.
    \item \textbf{src/features.py}: Computes rolling averages, robust scales numerical amounts, and produces naive and group-aware train/test splits.
    \item \textbf{src/sequence.py}: Frames structured tabular data into sliding temporal windows of length 6 to train recurrent layers.
    \item \textbf{src/models/rf\_rnn\_hybrid.py}: Manages Random Forest and SimpleRNN training, hyperparameter sweep over ensemble weights, and threshold tuning.
    \item \textbf{src/inference.py}: Wires together model predictions, executes explanations, and serves the scoring logic.
\end{itemize}

\subsection{Database or Data Design}
The data design defines the transactional feature contract. The dataset consists of 568,630 rows. Each row contains 31 raw columns: an incrementing `id`, 28 anonymized PCA features (`V1` to `V28`), a transaction `Amount` value, and a binary `Class` indicator (0 for legitimate, 1 for fraud). The feature engineering process expands this raw vector into 36 columns by generating:
\begin{itemize}
    \item \textbf{amount\_log1p}: Log-transformed amount to normalize extreme values.
    \item \textbf{amount\_robust\_scaled}: Amount centered by the median and scaled by the interquartile range (IQR).
    \item \textbf{time\_since\_last\_transaction\_proxy}: Elapsed time between successive events.
    \item \textbf{transaction\_frequency\_proxy}: Number of events within historical transaction windows.
    \item \textbf{rolling\_amount\_mean, rolling\_amount\_std}: Historical averages and deviations.
\end{itemize}

\section{Implementation Details}

\subsection{Machine Learning Implementation}

\subsubsection{Dataset Preparation \& Oversampling}
To handle extreme imbalance, \gls{smote} is applied to the minority class. Rather than oversampling the entire dataset beforehand, which leads to target leakage and overly optimistic test performance, \gls{smote} is applied strictly to the training fold. The test splits remain un-sampled, preserving the original class distribution to ensure representative evaluation metrics.

\subsubsection{Model Design}
The hybrid ensemble model consists of two components:
\begin{itemize}
    \item \textbf{Random Forest Component}: Trained on tabular feature vectors to capture static anomalies.
    \item \textbf{SimpleRNN Component}: Trained on sequences of historical data ($N \times 6 \times 36$) to learn temporal patterns. The network topology features:
    \begin{itemize}
        \item Input Layer: Shape $(6, 36)$
        \item SimpleRNN Layer: 64 hidden units, tanh activation
        \item Dense Layer: 1 unit, sigmoid activation (predicting fraud probability)
    \end{itemize}
\end{itemize}

\subsubsection{Training \& Ensemble Tuning}
The ensemble blends predictions using a weighted linear combination. The tuning script performs a grid search to find the optimal blending weights:
\begin{equation}
    Score = w_{\gls{rf}} \times P_{\gls{rf}} + w_{\gls{rnn}} \times P_{\gls{rnn}}
\end{equation}
subject to:
\begin{equation}
    w_{\gls{rf}} + w_{\gls{rnn}} = 1.0, \quad w_{\gls{rf}}, w_{\gls{rnn}} \ge 0
\end{equation}
The optimal threshold that maximizes the validation F1-score is chosen.

\section{Algorithms and Flowcharts}

The training and inference workflows are described in detail below.

\subsection{End-to-End Training Pipeline}
The end-to-end training pipeline orchestrates data loading, preprocessing, split creation, oversampling, and model sequence training:
\begin{enumerate}
    \item \textbf{Input}: Raw transactional dataset $\mathcal{D}_{raw}$.
    \item \textbf{Output}: Trained hybrid ensemble estimator $\{M_{rf}, M_{rnn}\}$, weights $\{w_{rf}, w_{rnn}\}$, optimal threshold $\theta$.
    \item \textbf{Data Cleaning}: Deduplicate and clean raw datasets.
    \item \textbf{Feature Engineering}: Apply Robust Scaling, log-amount transform, rolling frequency, and rolling averages.
    \item \textbf{Data Splitting}: Split into training and testing sets using a stratified, group-aware partitioning strategy to avoid leakages.
    \item \textbf{Oversampling}: Apply \gls{smote} strictly on training splits.
    \item \textbf{Tabular Training}: Fit the Random Forest component $M_{rf}$ on tabular features.
    \item \textbf{Temporal Windowing}: Re-shape training data into sequences of length 6.
    \item \textbf{Sequential Training}: Train the Recurrent Neural Network component $M_{rnn}$ on sequences.
    \item \textbf{Ensemble Blending}: Search optimal weight distribution $w_{rf}, w_{rnn}$ and threshold $\theta$ maximizing F1-score on validation sets.
\end{enumerate}

\subsection{Real-Time Inference Scoring Service}
The real-time scoring flow for online requests is executed as follows:
\begin{enumerate}
    \item \textbf{Input}: Transaction request vector $V_{in}$ containing 29 raw paper features.
    \item \textbf{Output}: Final classification verdict (FRAUD or LEGITIMATE) and \gls{shap} explainability.
    \item \textbf{Feature Mapping}: Transform raw 29-feature request into 36-feature vector $V_{model}$ using saved scaling and frequency parameters.
    \item \textbf{Static Scoring}: Feed $V_{model}$ into Random Forest to obtain probability $P_{rf}$.
    \item \textbf{Temporal Scoring}: Retrieve transaction history sequence and predict sequential probability $P_{rnn}$ using SimpleRNN.
    \item \textbf{Ensemble Blending}: Calculate final risk score: $Score = w_{rf} \times P_{rf} + w_{rnn} \times P_{rnn}$.
    \item \textbf{Threshold Gate}: If $Score \ge \theta$, label transaction as \textbf{FRAUD}; else label as \textbf{LEGITIMATE}.
    \item \textbf{SHAP Explanations}: Execute TreeSHAP explainer on $M_{rf}$ to extract local feature contributions.
\end{enumerate}
